%========================
% Table 1: eLife
%========================
\begin{table*}[h!]
\centering
\begin{tabular}{p{.95\textwidth}}
\textit{Technical} \\ \textexample{Whether complement dysregulation directly contributes to the pathogenesis of peripheral nervous system diseases, including sensory neuropathies, is unclear. We addressed this important question in a mouse model of ocular HSV-1 infection, where sensory nerve damage is a common clinical problem. Through genetic and pharmacologic targeting, we uncovered a central role for C3 in sensory nerve damage at the morphological and functional levels. Interestingly, CD4 T cells were central in facilitating this complement-mediated damage. This same C3/CD4 T cell axis triggered corneal sensory nerve damage in a mouse model of ocular graft-versus-host disease (GVHD). However, this was not the case in a T-dependent allergic eye disease (AED) model, suggesting that this inflammatory neuroimmune pathology is specific to certain disease etiologies. Collectively, these findings uncover a central role for complement in CD4 T cell-dependent corneal nerve damage in multiple disease settings and indicate the possibility for complement-targeted therapeutics to mitigate sensory neuropathies.

} \\
\textit{Lay} \\ \textexample{Most people have likely experienced the discomfort of an eyelash falling onto the surface of their eye. Or that gritty sensation when dust blows into the eye and irritates the surface. These sensations are warnings from sensory nerves in the cornea, the transparent tissue that covers the iris and pupil. Corneal nerves help regulate blinking, and control production of the tear fluid that protects and lubricates the eye. But if the cornea suffers damage or infection, it can become inflamed. Long-lasting inflammation can damage the corneal nerves, leading to pain and vision loss. If scientists can identify how this happens, they may ultimately be able to prevent it. To this end, Royer et al. have used mice to study three causes of hard-to-treat corneal inflammation. The first is infection with herpes simplex virus (HSV-1), which also causes cold sores. The second is eye allergy, where the immune system overreacts to substances like pollen or pet dander. And the third is graft-versus-host disease (GVHD), an immune disorder that can affect people who receive a bone marrow transplant. Royer et al. showed that HSV-1 infection and GVHD – but not allergies – made the mouse cornea less sensitive to touch. Consistent with this, microscopy revealed damage to corneal nerves in the mice with HSV-1 infection and those with GVHD. [...]

} \\
\end{tabular}
\caption{Example technical and lay summaries from eLife.}
\label{tab:dataset-examples-elife}
\end{table*}

%========================
% Table 2: PLOS
%========================
\begin{table*}[h!]
\centering
\begin{tabular}{p{.95\textwidth}}
\textit{Technical} \\
\textexample{The 3-O-sulfotransferase (3-OST) family catalyzes rare modifications of glycosaminoglycan chains on heparan sulfate proteoglycans, yet their biological functions are largely unknown. Knockdown of 3-OST-7 in zebrafish uncouples cardiac ventricular contraction from normal calcium cycling and electrophysiology by reducing tropomyosin4 (tpm4) expression. Normal 3-OST-7 activity prevents the expansion of BMP signaling into ventricular myocytes, and ectopic activation of BMP mimics the ventricular noncontraction phenotype seen in 3-OST-7 depleted embryos. In 3-OST-7 morphants, ventricular contraction can be rescued by overexpression of tropomyosin tpm4 but not by troponin tnnt2, indicating that tpm4 serves as a lynchpin for ventricular sarcomere organization downstream of 3-OST-7. Contraction can be rescued by expression of 3-OST-7 in endocardium, or by genetic loss of bmp4. Strikingly, BMP misregulation seen in 3-OST-7 morphants also occurs in multiple cardiac noncontraction models, including potassium voltage-gated channel gene kcnh2 affected in Romano-Ward syndrome and long-QT syndrome. [...]} \\
\textit{Lay} \\
\textexample{A highly complex environment at the cell surface and in the space between cells is thought to modulate cell behavior. Heparan sulfate proteoglycans are cell surface and extracellular matrix molecules that are covalently linked to long chains of repeating sugar units called glycosaminoglycan chains. These chains can be subjected to rare modifications and they are believed to influence specific cell signaling events in a lineage-specific fashion in what is called the “glycocode.” Here we explore the functions of one member of a family of enzymes, 3-O-sulfotransferases (3-OSTs), that catalyzes a rare modification (3-O-sulfation) of glycosaminoglycans in zebrafish. We show that knockdown of 3-OST-7 results in a very specific phenotype, including loss of cardiac ventricle contraction. Knockdown of other 3-OST family members did not result in the same phenotype, suggesting that distinct 3-OST family members have distinct functions in vertebrates and lending in vivo evidence for the glycocode hypothesis. [...]} \\
\end{tabular}
\caption{Example technical and lay summaries from PLOS.}
\label{tab:dataset-examples-plos}
\end{table*}

%========================
% Table 3: Eureka
%========================
\begin{table*}[h!]
\centering
\begin{tabular}{p{.95\textwidth}}
\textit{Technical} \\ \textexample{Extracellular vesicles (EVs) are small vesicles released by cells to aid cell–cell communication and tissue homeostasis. Human islet amyloid polypeptide (IAPP) is the major component of amyloid deposits found in pancreatic islets of patients with type 2 diabetes (T2D). IAPP is secreted in conjunction with insulin from pancreatic cells to regulate glucose metabolism. Here, using a combination of analytical and biophysical methods in vitro, we tested whether EVs isolated from pancreatic islets of healthy patients and patients with T2D modulate IAPP amyloid formation. We discovered that pancreatic EVs from healthy patients reduce IAPP amyloid formation by peptide scavenging, but T2D pancreatic and human serum EVs have no effect. In accordance with these differential effects, the insulin:C-peptide ratio and lipid composition differ between EVs from healthy pancreas and EVs from T2D pancreas and serum. It appears that healthy pancreatic EVs limit IAPP amyloid formation via direct binding as a tissue-specific control mechanism.} \\
\textit{Lay} \\ \textexample{A step closer to a cure for adult-onset diabetes. Professor and head of the Chemical Biology division at the Department of Biology and Biological Engineering, she leads a research team focusing on metalloproteins and proteins that fold incorrectly. Exosomes in patients with the disease don’t have the same ability. This discovery, by a research collaboration between Chalmers University of Technology and AstraZeneca, takes us a step closer to a cure for type 2 diabetes. Proteins are the body’s workhorses, carrying out all the tasks in our cells. A protein is a long chain of amino acids that must be folded into a specific three-dimensional structure to function properly. Sometimes, however, they misbehave and aggregate—clump together—into long fibers called amyloids, which can cause diseases. It was previously known that type 2 diabetes is caused by a protein aggregating in the pancreas. “What we’ve found is that exosomes secreted by the cells in the pancreas stop that process in healthy people and protect them from type 2 diabetes, while the exosomes of diabetes patients do not,” says Professor Pernilla Wittung-Stafshede, who headed the study. [...]} \\
\end{tabular}
\caption{Example technical and lay summaries from Eureka.}
\label{tab:dataset-examples-eureka}
\end{table*}

%========================
% Table 4: SciNews
%========================
\begin{table*}[h!]
\centering
\begin{tabular}{p{.95\textwidth}}
\textit{Technical} \\ \textexample{Observational studies suggest that lower educational attainment (EA) may be associated with risky alcohol use behaviors; however, these findings may be biased by confounding and reverse causality. We performed two-sample Mendelian randomization (MR) using summary statistics from recent genome-wide association studies (GWAS) with $<$780,000 participants to assess the causal effects of EA on alcohol use behaviors and alcohol dependence (AD). Fifty-three independent genome-wide significant SNPs previously associated with EA were tested for association with alcohol use behaviors. We show that while genetic instruments associated with increased EA are not associated with total amount of weekly drinks, they are associated with reduced frequency of binge drinking $\geq$6 drinks ($\beta_\mathrm{IVW} = -0.198$, 95\% CI, $-0.297$ to $-0.099$, $P_\mathrm{IVW} = 9.14 \times 10^{-5}$) [...]} \\
\textit{Lay} \\ \textexample{Higher educational attainment—spending more years in education—may impact people's drinking behaviour and reduce their risk of alcohol dependence, according to a study published in Molecular Psychiatry. Alcohol consumption is a major risk factor for death and disability worldwide. Identifying factors associated with how much, how often and what people drink may be important for developing and improving intervention and treatment strategies. Previous studies have suggested that educational attainment may influence drinking, but with conflicting results. To assess the possible effects of educational attainment on alcohol use behaviours and alcohol dependence, a team of researchers at the National Institutes of Health, USA used two-sample Mendelian randomisation statistical methods. [...] Dr. Lohoff said: "It is important to understand that while these genetic variants allow us to investigate the possible effect of educational attainment on alcohol consumption and alcohol dependence, this doesn't mean that educational attainment can't be modified. The possible effect of educational attainment on drinking that we show in this study suggests that increasing educational attainment may be a useful target for prevention programs against problematic alcohol use, alcohol dependence, and their consequences." [...]} \\
\end{tabular}
\caption{Example technical and lay summaries from SciNews.}
\label{tab:dataset-examples-scinews}
\end{table*}
